\documentclass{standalone}
\usepackage{xparse}
\usepackage{tikz}
\usepackage{graphicx}
\usetikzlibrary{shapes}

\newcount\xhexes
% The hex size is determined by the number of hexes you want horizontally
% Change the number below to modify the hexes in the x axis
\xhexes=10

% Use makeatletter so we can use strip@pt
\makeatletter
\newcommand*{\getlength}[1]{\strip@pt#1}
\makeatother

\newdimen\imagewidth
\imagewidth=\textwidth

\newdimen\imageheight
\settoheight{\imageheight}{%
  \includegraphics[width=\imagewidth,keepaspectratio]{example-image}%
}

\newdimen\hexdim
\hexdim=\numexpr\getlength{\imagewidth}/(\xhexes) pt\relax

\imageheight=258pt
\newcount\yhexes
\yhexes=\numexpr\getlength{\imageheight} / \getlength{\hexdim}\relax

\begin{document}

  \the\imageheight
  \the\hexdim
  \the\xhexes

  \begin{tikzpicture} [hexa/.style= {shape=regular polygon,
                                     regular polygon sides=6,
                                     minimum size=\hexdim, draw,
                                     inner sep=0,anchor=south,
                                     opacity=.5}]

      \node[anchor=south west,inner sep=0] (image) at (0,0) {\includegraphics[width=\imagewidth]{example-image}};
      \foreach \j in {0,...,\xhexes}{% 
           \ifodd\j 
               \foreach \i in {0,...,\yhexes}{\node[hexa] (h\j;\i) at ({\j*\hexdim/2+\j*\hexdim/4+\hexdim/2},{(\i*\hexdim+\hexdim/2)*sin(60)}) {\j;\i};}        
          \else
               \foreach \i in {0,...,\yhexes}{\node[hexa] (h\j;\i) at ({\j*\hexdim/2+\j*\hexdim/4+\hexdim/2},{\i*\hexdim*sin(60)}) {\j;\i};}
          \fi}
  \end{tikzpicture}

\end{document}      
